% BHU PYQ -- Professional Exam Template
\documentclass[11pt,a4paper]{article}

\usepackage[a4paper,top=2.5cm,bottom=2.5cm,left=2.8cm,right=2.8cm,headheight=14pt]{geometry}
\usepackage{amsmath,amssymb}
\usepackage[T1]{fontenc}
\usepackage[utf8]{inputenc}
\usepackage{lmodern}
\usepackage{microtype}
\usepackage{fancyhdr}
\usepackage{enumitem}
\usepackage{hyperref}
\usepackage{lastpage}
\usepackage{parskip}
\usepackage{tikz}
\usetikzlibrary{arrows.meta,shapes,positioning}

\hypersetup{colorlinks=true,urlcolor=black,linkcolor=black,
  pdftitle={B.Sc. (Hons.) Zoology Sem-I ZOB-101 (Old Course 2012-13) -- BHU PYQ}}

\pagestyle{fancy}
\fancyhf{}
\renewcommand{\headrulewidth}{0.4pt}
\renewcommand{\footrulewidth}{0.4pt}
\fancyhead[L]{\small   \textit{Banaras Hindu University}}
\fancyhead[R]{\small   \textit{Zoology Paper: ZOB-101 (Old 2012)}}
\fancyfoot[C]{\small Page  \thepage\ of \pageref{LastPage}}


\newlist{parts}{enumerate}{3}
\setlist[parts,1]{label=(\alph*),leftmargin=*,itemsep=4pt,topsep=3pt}
\setlist[parts,2]{label=(\roman*),leftmargin=*,itemsep=2pt,topsep=2pt}
\setlist[parts,3]{label=(\arabic*),leftmargin=*,itemsep=2pt,topsep=2pt}

\newcommand{\pts}[1]{\hfill   \textit{[#1]}}


\begin{document}


\begin{center}
  {\large\bfseries BANARAS HINDU UNIVERSITY}\\[2pt]
  {
\normalsize B.Sc. (Hons.) Semester I Examination, 2016-17}\\[6pt]
  \rule{0.8\linewidth}{0.5pt}\\[6pt]
  {\large\bfseries Zoology}\\[4pt]
  {\large\bfseries Paper: ZOB-101 --- Animal Diversity \& Fundamentals of Cell Biology}\\[2pt]
  {
\normalsize   \textbf{(Old Course 2012-13 Syllabus)}}\\[4pt]
  {
\normalsize Time: Three hours \hfill Full Marks: 70}
\end{center}

\rule{\linewidth}{0.4pt}

\smallskip



\noindent   \textit{Note: Answer three questions from Section-A and three questions from Section-B. Question No. 1 in both sections is compulsory. Use a separate answer book for each Section. The figures in the right-hand margin indicate marks.}

\bigskip

\begin{center}
   \textbf{SECTION -- A} \\
   \textbf{(Animal Diversity) (Marks: 35)}
\end{center}

\smallskip



\noindent   \textit{Note: Answer three questions including Question No. 1 which is compulsory.}

\medskip


\noindent   \textbf{Question 1.} Answer the following parts:

\begin{parts}
  \item State whether the following statements are True or False. Justify your answer in 2--3 sentences:\pts{$3  \times 2 = 6$}
  
\begin{parts}
    \item Protozoans belong to subphylum Cnidospora and Sporozoa produce same kind of spores.
    \item The secondary host of   \textit{Fasciola hepatica} is sheep.
    \item The lung fish belong to subclass Sarcopterygii.
  \end{parts}

  \item Differentiate between the following:\pts{$2  \times 2 = 4$}
  
\begin{parts}
    \item Homology and analogy
    \item Flightless and flying birds.
  \end{parts}

  \item Define metamerism.\pts{1}
\end{parts}

\medskip


\noindent   \textbf{Question 2.} Answer the following parts:\pts{$7 + 5 = 12$}

\begin{parts}
  \item Classify phylum Echinodermata up to classes giving characteristic features and examples of each class.
  \item Describe different types of scales in fishes.
\end{parts}

\medskip


\noindent   \textbf{Question 3.} Answer the following parts:\pts{$7 + 5 = 12$}

\begin{parts}
  \item Give an account of salient features of phylum Porifera. Classify up to classes giving characteristic features and examples of each class.
  \item Describe the canal system in Phylum Porifera.
\end{parts}

\medskip


\noindent   \textbf{Question 4.} Write short notes on any three of the following:\pts{$3  \times 4 = 12$}

\begin{parts}
  \item Air sacs in birds
  \item Protochordates
  \item Chiroptera
  \item Sarcomastigophora.
\end{parts}

\pagebreak

\begin{center}
   \textbf{SECTION -- B} \\
   \textbf{(Fundamentals of Cell Biology) (Marks: 35)}
\end{center}

\smallskip



\noindent   \textit{Note: Answer three questions including Question No. 1 which is compulsory.}

\medskip


\noindent   \textbf{Question 1.} Answer the following parts:

\begin{parts}
  \item Write whether the following statements are True or False, giving reasons in 2--3 sentences:\pts{$2  \times 4 = 8$}
  
\begin{parts}
    \item Eukaryotic cells have cell wall made-up of peptidoglycans.
    \item Cellular respiration occurs on the inner layer of plasma membrane of alveolar cells only.
    \item Polytene chromosomes are interphase chromosomes.
    \item Homozygous loss of function of an oncogene leads to the transformation of cells from normal to cancerous state.
  \end{parts}

  \item Define the following:\pts{$1  \times 3 = 3$}
  
\begin{parts}
    \item Prokaryote
    \item Cell cycle
    \item Antibody.
  \end{parts}
\end{parts}

\medskip


\noindent   \textbf{Question 2.} Write notes on the following:\pts{$4  \times 3 = 12$ (or $3  \times 4 = 12$)}

\begin{parts}
  \item Function of rough endoplasmic reticulum
  \item Image formation in a student's compound microscope
  \item Significance of meiosis.
\end{parts}

\medskip


\noindent   \textbf{Question 3.} Answer the following parts:\pts{$6 + 6 = 12$}

\begin{parts}
  \item Draw a labeled flow chart diagram to depict biogenesis of ribosomes (no descriptions required).
  \item Illustrate how the sorting of lysosomal proteins takes place and discuss the function of a lysosome.
\end{parts}

\medskip


\noindent   \textbf{Question 4.} Answer the following parts:\pts{$6 + 6 = 12$}

\begin{parts}
  \item Differentiate between mitosis and meiosis showing clearly the changes in DNA content at each phase of division.
  \item Illustrate structure and function of peroxisome.
\end{parts}

\vfill
\rule{\linewidth}{0.4pt}

\begin{center}\small
\href{https://bhu-exam.vercel.app/}{Click Here}\\[4pt]\href{https://me-aryan.vercel.app/}{Click Here}\\[4pt]\href{https://quashan-me.vercel.app/}{Click Here}
\end{center}

\end{document}
